% items/figures/architecture.tex
% Архитектурная схема для §3 ВКР.
% Требует: \usepackage{tikz} и tikzlibrary: positioning, arrows.meta, fit, calc, backgrounds

\begin{figure}[htbp]
\centering

% Локальный хелпер: «§ 4.3» как кликабельная ссылка на нужный label
\providecommand{\secref}[1]{\hyperref[#1]{\textsection\,\ref*{#1}}}

\begin{tikzpicture}[
  >=Stealth,
  font=\footnotesize,
  every node/.style={font=\footnotesize, align=left},
  % Большой блок-слой
  layer/.style={
    draw=black!55,
    rounded corners=1.5pt,
    inner xsep=4mm,
    inner ysep=2.5mm,
  },
  % Внутренний модуль (в Слое 3)
  module/.style={
    draw=black!40,
    rounded corners=1pt,
    inner xsep=2mm,
    inner ysep=2mm,
    text width=44mm,
  },
  % Стрелка между слоями
  arr/.style={->, thick, black!55},
  % Подпись на стрелке
  arrlabel/.style={font=\scriptsize\itshape, text=black!60, midway, right=1mm},
  % Заголовок слоя
  ltitle/.style={font=\footnotesize\bfseries},
]

% ── Контекст пользователя ────────────────────────────────────────────────
\node[layer, text width=150mm, fill=black!3] (ctx) {
  {\footnotesize\bfseries Контекст пользователя}\\[1mm]
  Тип пользователя \;\textperiodcentered\;
  Уровень образования \;\textperiodcentered\;
  Уровень погружённости \;\textperiodcentered\;
  Цель просмотра
};

% ── Слой 1 ──────────────────────────────────────────────────────────────
\node[layer, text width=150mm, below=8mm of ctx] (l1) {
  {\footnotesize\bfseries Сбор данных}\\[1mm]
  YouTube URL $\to$ метаданные \;\textperiodcentered\;
  аудиопоток \;\textperiodcentered\;
  видеокадры \;\textperiodcentered\; транскрипт
};

% ── Слой 2 ──────────────────────────────────────────────────────────────
\node[layer, text width=150mm, below=4mm of l1] (l2) {
  {\footnotesize\bfseries Извлечение признаков}\\[1mm]
  аудио-признаки \quad|\quad
  видео-признаки \quad|\quad
  текст-признаки
};

% ── Слой 3: три модуля + общая рамка ────────────────────────────────────
% Сначала ставим модули, потом дорисовываем рамку «под ними» через background layer
\node[module, below=10mm of l2.south west, xshift=4mm, anchor=north west] (m31) {
  \textbf{Техническое качество}\\[0.8mm]
  аудио \;\secref{sec:method-audio}\\
  видео \;\secref{sec:method-video}
};
\node[module, right=2mm of m31, text width=58mm] (m32) {
  \textbf{Флаги восприятия} \;\secref{sec:method-flags}\\[0.8mm]
  Аудио дисбаланс\\
  Смысловая унылость\\
  Эмоциональный дисбаланс\\
  Визуальный шум
};
\node[module, right=2mm of m32] (m33) {
  \textbf{Профиль подачи} \;\secref{sec:method-profile}\\[0.8mm]
  Две независимые оси:\\
  академичность \;$\times$\; инструментальность
};
\node[ltitle, above=1.5mm of m31.north west, anchor=south west] (l3title) {Слой 3 — анализ};

\begin{scope}[on background layer]
  \node[layer, fit=(l3title)(m31)(m32)(m33)] (l3) {};
\end{scope}

% ── Слой 4 ──────────────────────────────────────────────────────────────
\node[layer, text width=150mm, below=4mm of l3] (l4) {
  {\footnotesize\bfseries Синтез и выдача}\\[1mm]
  Нарративная суммаризация \;\secref{sec:method-narrative}
  \quad+\quad
  Метрика соответствия запросу \;\secref{sec:method-relevance}\\[1mm]
  {\itshape выдача:}
  техническое качество \;\textperiodcentered\;
  флаги \;\textperiodcentered\;
  профиль подачи \;\textperiodcentered\;
  нарратив с таймкодами \;\textperiodcentered\;
  интегральная оценка
};

% ── Стрелки ─────────────────────────────────────────────────────────────
\draw[arr] (ctx) -- node[arrlabel] {влияет на все нижеследующие слои} (l1);
\draw[arr] (l1) -- (l2);
\draw[arr] (l2.south) -- (l3.north);
\draw[arr] (l3.south) -- (l4.north);

\end{tikzpicture}
\caption{Архитектура системы: четырёхслойный каскад с пользовательским
контекстом, пронизывающим интерпретацию на всех уровнях.
Ссылки указывают на параграфы \cref{sec:methodology}, раскрывающие модули.}
\label{fig:architecture}
\end{figure}